% !TeX root = ../navier-stokes-manuscript.tex
\subsection{The finite piece that reaches a time slice}\label{sub:BodyFiniteAdjacentRectangles}

Let \(I_*=(0,T_*)\) denote the interval on a proposed finite maximal branch,
and first stop at a strict cutoff \(T<T_*\).  Control of \(v\) one spatial
step above \(H^M\) has no reason by itself to arrive continuously at \(t=T\).
Its time derivative must be controlled one spatial step below \(H^M\) as well.

Let \(A=(1-\Delta)^{1/2}\).  For a smooth Hilbert-valued function \(v\),
\begin{equation}\label{eq:BodyAdjacentTraceIdentity}
  \norm{v(b)}_{H^M}^2-
  \norm{v(a)}_{H^M}^2
  =2\operatorname{Re}\int_a^b
  \pair{A^{M-1}v_t(t)}{A^{M+1}v(t)}_{L^2}\,dt.
\end{equation}
Cauchy--Schwarz places the middle trace between these two adjacent spatial
controls
\begin{equation}\label{eq:BodyAdjacentTracePair}
  v\in L^2(0,T;H^{M+1}),
  \qquad
  v_t\in L^2(0,T;H^{M-1}),
\end{equation}
and the conclusion is
\begin{equation}\label{eq:BodyAdjacentTraceConclusion}
  v\in C([0,T];H^M);
\end{equation}
the approximation argument and strong representative are proved in
\Cref{lem:AdjacentHilbertScaleTrace}.

At \(M=2\), one concrete adjacent pair is
\begin{equation}\label{eq:FirstConcreteAdjacentPair}
  u\in L^2(0,T;H^3),
  \qquad
  u_t\in L^2(0,T;H^1)
  \quad\Longrightarrow\quad
  u\in C([0,T];H^2).
\end{equation}
The velocity lies one spatial step above the desired trace and its next
temporal response lies one spatial step below it.  Their pairing reaches the
middle.

Now apply that trace to \(V_n\), whose derivative is \(V_{n+1}\).  To retain
the responses through order \(N\) at a common height \(M\), collect
\begin{equation}\label{eq:BodyRectangleNorm}
\begin{aligned}
  \mathfrak R_{N,M}(u;T)
  :={}&
  \sum_{n=0}^{N}
  \norm{V_n}_{L^2(0,T;H^{M+1})}^2
  \\
  &+
  \sum_{n=0}^{N}
  \norm{V_{n+1}}_{L^2(0,T;H^{M-1})}^2.
\end{aligned}
\end{equation}
When \(M=0\), the lower edge is read in \(H^{-1}\).  The temporal depth
\(N\) and spatial height \(M\) remain independent choices.
The number \(\mathfrak R_{N,M}(u;T)\) measures the size of this finite view of
the history.

\Needspace{12\baselineskip}
The \((N,M)=(1,2)\) example contains two adjacent traces sharing \(V_1\):
\begin{center}
\captionsetup{hypcap=false}
\captionof{table}{The \((1,2)\) rectangle.  The shared response closes the first trace and begins the second.}
\label{tab:FirstAdjacentRectangle}
\renewcommand{\arraystretch}{1.3}
\begin{tabular}{@{}c c c@{}}
\toprule
response & above the trace & below the trace \\
\midrule
\(V_0=u\) & \(V_0\in L^2H^3\) & \(V_1\in L^2H^1\) \\
\(V_1=u_t\) & \(V_1\in L^2H^3\) & \(V_2\in L^2H^1\) \\
\bottomrule
\end{tabular}
\end{center}

The norm alone does not remember which field produced these controls.  Write
\(\cR_{N,M}[u_0;T]\) for the velocity block together with the pressure
responses reconstructed by \eqref{eq:BodyTemporalPressureRecurrence}.  Every
entry is a derivative of the same velocity or the normalized pressure forced
by that velocity.

Classical smoothness on a strict subinterval makes each chosen rectangle
finite:
\begin{equation}\label{eq:StrictSubintervalRectangleBound}
  \mathfrak R_{N,M}(u;T)<\infty
  \qquad(T<T_*,\;N,M<\infty).
\end{equation}
But the terminal question changes the order of the quantifiers:
\[
  \forall\,T<T_*:\quad \mathfrak R_{N,M}(u;T)<\infty
  \quad\not\Longrightarrow\quad
  \sup_{0<T<T_*}\mathfrak R_{N,M}(u;T)<\infty.
\]
Fixing \((N,M)\) first leaves one demand: a bound that survives every cutoff
approaching \(T_*\).  This is the whole-terminal rectangle estimate.
